\section{Empirical Evaluation}
\label{section:evaluation}
To rigorously evaluate the effectiveness of our proposed RL-enhanced IndexTTS2 framework in recovering and controlling emotions for low-resource languages, we employ a combination of objective and subjective evaluations.

\subsection{Performance Metrics}
We plan to adopt several key performance metrics to comprehensively assess speech intelligibility, speaker fidelity, and emotional expressiveness.
\begin{itemize}
    \item \textbf{Objective Metrics}:
    \begin{itemize}
        \item \textbf{Word Error Rate (WER)}: Evaluates speech intelligibility and articulation accuracy using ASR models like FunASR or Whisper \citep{Radford2022RobustSR}.
        \item \textbf{Speaker Similarity (SS)}: Computed as the cosine similarity between speaker embeddings extracted from a pretrained speaker recognition model to ensure the timbre is maintained after RL fine-tuning.
        \item \textbf{Emotion Similarity (ES) / Accuracy}: Measured using the open-source emotion2vec model to objectively verify if the generated speech matches the target emotion category.
    \end{itemize}
    \item \textbf{Subjective Metrics (MOS Framework)}:
    \begin{itemize}
        \item \textbf{Multidimensional MOS}: We will conduct human evaluations based on a 1-5 scale for Quality MOS (QMOS), Similarity MOS (SMOS), Prosody MOS (PMOS), and Emotion MOS (EMOS).
        \item \textbf{Emotion Intensity Test (EIT)}: A pairwise comparison task where listeners identify the stronger emotional sample among weak, medium, and strong synthesized speech. 
        \item \textbf{Emphasis Accuracy Test (EAT)}: Assesses the consistency between the model's predicted emphasis positions (local prosody) and those perceived by human listeners.
    \end{itemize}
\end{itemize}
The combination of EIT and EAT is crucial as our RL method explicitly optimizes for fine-grained intensity and local emphasis, which standard MOS cannot fully capture.

\subsection{Baseline Methods}
We compare our method against several state-of-the-art zero-shot and emotionally controllable TTS models.
\begin{itemize}
    \item \textbf{CosyVoice2} \citep{Du2024CosyVoice2S}: A highly scalable, LLM-based zero-shot text-to-speech synthesizer that supports predefined emotion instructions.
    % \item \textbf{MaskGCT} \citep{Wang2024MaskGCTZT}: A zero-shot text-to-speech model utilizing a masked generative codec transformer for non-autoregressive synthesis.
    % \item \textbf{F5-TTS} \citep{Chen2024F5TTSAF}: A flow-matching based non-autoregressive TTS model known for fluent and faithful speech generation.
    % \item \textbf{Spark-TTS} \citep{Wang2025SparkTTSAE}: An efficient LLM-based TTS model that utilizes single-stream decoupled speech tokens.
    % \item \textbf{EmoSphere++} \citep{Cho2024EmoSphereEZ}: An emotion-controllable zero-shot TTS baseline that models emotions via an emotion-adaptive spherical vector.
    \item \textbf{IndexTTS / IndexTTS2} \citep{Deng2025IndexTTSAI, Zhou2025IndexTTS2AB}: The foundational autoregressive models of our work, providing baseline performance for duration control and semantic generation before RL enhancement.
\end{itemize}

\subsection{Benchmark Tasks or Datasets}
To simulate the low-resource fine-tuning scenario and evaluate the catastrophic forgetting of emotions, we utilize the following datasets:
\begin{itemize}
    \item \textbf{Base Pre-training Datasets}: The model is initially pre-trained on large-scale datasets such as the Emilia dataset and LibriSpeech to establish robust zero-shot capabilities.
    \item \textbf{Emotional Fine-tuning Datasets}: We utilize the Emotional Speech Database (ESD) and the Expresso dataset, which provide rich annotations for discrete emotions, continuous intensities, and emphasis. 
    \item \textbf{Low-Resource Simulation Task}: We will construct a specific low-resource task by taking a subset (e.g., 1 to 5 hours) of a target low-resource language corpus. The model will be fine-tuned on this subset, followed by our RL optimization stage, to measure its ability to retain and express the emotional traits learned from the ESD and Expresso datasets.
    \item \textbf{Standard Evaluation Sets}: For standard synthesis quality and stability testing, we will report metrics on SeedTTS test-zh, SeedTTS test-en, and LibriSpeech-test-clean.
\end{itemize}